%; whizzy chapter -dvi
% -initex iniptex -latex platex -format platex -bibtex jbibtex -fmt fmt
% 以上 whizzytex を使用する場合の設定。
 
%     Tokyo Debian Meeting resources
%     Copyright (C) 2012 Junichi Uekawa
%     Copyright (C) 2011, 2015 Nobuhiro Iwamatsu

%     This program is free software; you can redistribute it and/or modify
%     it under the terms of the GNU General Public License as published by
%     the Free Software Foundation; either version 2 of the License, or
%     (at your option) any later version.

%     This program is distributed in the hope that it will be useful,
%     but WITHOUT ANY WARRANTY; without even the implied warranty of
%     MERCHANTABILITY or FITNESS FOR A PARTICULAR PURPOSE.  See the
%     GNU General Public License for more details.

%     You should have received a copy of the GNU General Public License
%     along with this program; if not, write to the Free Software
%     Foundation, Inc., 51 Franklin St, Fifth Floor, Boston, MA  02110-1301 USA

%  preview (shell-command (concat "evince " (replace-regexp-in-string "tex$" "pdf"(buffer-file-name)) "&"))

%%ここからヘッダ開始。

\documentclass[mingoth,a4paper]{jsarticle}
\usepackage{monthlyreport}
% 日付を定義する、毎月変わります。
\newcommand{\debmtgyear}{2015}
\newcommand{\debmtgmonth}{11}
\newcommand{\debmtgdate}{21}
% started from zero:
% (let ((year 2013) (month 7)) (+ (* (- year 2005) 12) month -1))
\newcommand{\debmtgnumber}{133}

\begin{document}

\begin{titlepage}
\thispagestyle{empty}
% タイトルページ:編集必要な部分は最初のマクロに飛ばすこと

\vspace*{-2cm}
第\debmtgnumber{}回 東京エリア Debian 勉強会資料\\
\hspace*{-2cm}
\includegraphics{image2012-natsu/dotdeb.pdf}\\
\hfill{}\debmtgyear{}年\debmtgmonth{}月\debmtgdate{}日

% ここはアップデートすること
% 全角文字にしないとフォントのサイズが合わないので注意
\rotatebox{10}{\fontsize{30}{30} {\gt 特集：kFreeBSD セットアップガイド}}\\

\vspace*{-2cm}
\hfill{}\includegraphics[height=6cm]{image200502/openlogo-nd.eps}
\end{titlepage}

\newpage

\begin{minipage}[b]{0.2\hsize}
 \definecolor{titleback}{gray}{0.9}
 \colorbox{titleback}{\rotatebox{90}{\fontsize{80}{80} {\gt デビアン勉強会} }}
\end{minipage}
\begin{minipage}[b]{0.8\hsize}
\hrule
\vspace{2mm}
\hrule
\begin{multicols}{2}
\tableofcontents
\end{multicols}
\vspace{2mm}
\hrule
\end{minipage}

\dancersection{事前課題}{野島 貴英}

今回の事前課題は以下です:
\begin{enumerate}
\item hack time に何をしますか？
\item (任意回答) 本勉強会をどこでお知りになりましたか？
\end{enumerate}
この課題に対して提出いただいた内容は以下です。
\begin{multicols}{2}
{\small
\begin{prework}{ wskoka }
  \begin{enumerate}
  \item Q.hack time $B$K2?$r$7$^$9$+!)(B\\
    A. Tilera$B$N(BDebian$B:n$j$?$$!#(B
  \item Q.$BK\JY6/2q$r$I$3$G$*CN$j$K$J$j$^$7$?$+!)(B\\
    A. dots$B$N%5%$%HWQWKCf$KH/8+!#(B
  \end{enumerate}
\end{prework}

\begin{prework}{ henrich }
  \begin{enumerate}
  \item Q.hack time $B$K2?$r$7$^$9$+!)(B\\
    A. $BK]Lu:n6H$r?J$a$?$$$H;W$$$^$9!#(B
  \end{enumerate}
\end{prework}

\begin{prework}{ hiroyuki.nagata }
  \begin{enumerate}
  \item Q.hack time $B$K2?$r$7$^$9$+!)(B\\
    A. $B$J$K$+%W%m%0%i%`$r=q$/(B
  \item Q.$BK\JY6/2q$r$I$3$G$*CN$j$K$J$j$^$7$?$+!)(B\\
    A. $B$?$7$+(BTwitter
  \end{enumerate}
\end{prework}

\begin{prework}{ dictoss }
  \begin{enumerate}
  \item Q.hack time $B$K2?$r$7$^$9$+!)(B\\
    A. Debian GNU/kFreeBSD$B$NF0:n%F%9%H(B
  \item Q.$BK\JY6/2q$r$I$3$G$*CN$j$K$J$j$^$7$?$+!)(B\\
    A. $B8}%3%_(B
  \end{enumerate}
\end{prework}

\begin{prework}{ rosh }
  \begin{enumerate}
  \item Q.hack time $B$K2?$r$7$^$9$+!)(B\\
    A. BTS$B$N:n6H(B
  \end{enumerate}
\end{prework}

\begin{prework}{ koedoyoshida }
  \begin{enumerate}
  \item Q.hack time $B$K2?$r$7$^$9$+!)(B\\
    A. $B$"$s$I$-$e$a$s$F$C$I$G$S$"$sE_9f$NJT=8(B
  \item Q.$BK\JY6/2q$r$I$3$G$*CN$j$K$J$j$^$7$?$+!)(B\\
    A. $BA02s$NJY6/2q(B
  \end{enumerate}
\end{prework}

\begin{prework}{ yy\_y\_ja\_jp }
  \begin{enumerate}
  \item Q.hack time $B$K2?$r$7$^$9$+!)(B\\
    A. DDTSS
  \end{enumerate}
\end{prework}

\begin{prework}{ $B$+$8(B }
  \begin{enumerate}
  \item Q.hack time $B$K2?$r$7$^$9$+!)(B\\
    A. XCB$B%W%m%0%i%_%s%0(B
  \end{enumerate}
\end{prework}

\begin{prework}{ $BLnEg(B }
  \begin{enumerate}
  \item Q.hack time $B$K2?$r$7$^$9$+!)(B\\
    A. DDTSS$B$H$+(B
  \item Q.$BK\JY6/2q$r$I$3$G$*CN$j$K$J$j$^$7$?$+!)(B\\
    A. $B44;v$@$C$?$j$G$7$F(B...
  \end{enumerate}
\end{prework}




}
\end{multicols}


\dancersection{Debian Trivia Quiz}{野島 貴英}

Debianの昨今の話題についてのQuizです。

今回の出題範囲は\url{debian-devel-announce@lists.debian.org} や \url{debian-news@lists.debian.org}に投稿された
内容などからです。

\begin{multicols}{2}
%; whizzy-master ../debianmeetingresume201311.tex
% $B0J>e$N@_Dj$r$7$F$$$k$?$a!"$3$N%U%!%$%k$G(B M-x whizzytex $B$9$k$H!"(Bwhizzytex$B$,MxMQ$G$-$^$9!#(B
%

\santaku
{2015/10/22$B$N(BDPN$B$N%a!<%k$+$iDj4|E*$K$J$,$l$F$$$?$$$/$D$+$N%H%T%C%/$,(BWeb$B$N$_$K7G<($5$l$k$h$&$K$J$j$^$7$?!#$I$3$K7G<($5$l$k!)(B}
{http://www.debian.or.jp/}
{http://www.debian.org/}
{http://bits.debian.org/}
{C}
{2015/10/22$B$N(BDPN$B$N%a!<%k$+$i!"=>MhN.$l$F$$$?(BDPN$B$N%a!<%k$N9=@.$,:~?7$5$l$^$7$?!#$3$A$i$KH<$$!"%;%-%e%j%F%#$K$D$$$F$N%"%J%&%s%9$H!"?7$7$$(BDD/DM$B$N%"%J%&%s%9$O!"(BDPN$B$N%a!<%k$K$O4^$^$l$:!"(BWeb$B%Z!<%8$K7G:\$5$l$k$N$_$H$J$j$^$7$?!#(B}

\santaku
{2015/11/7$B$K$F!"$H$"$k%$%s%9%?%s%H%a%C%;%s%8%c!<MQ%W%m%H%3%k$N%5!<%S%9$,A4(BDebian Developer$B$G;H$($k$h$&$K$J$C$?$H$N%"%J%&%s%9$,$"$j$^$7$?!#%W%m%H%3%k$NL>A0$O<!$N$I$l!)(B}
{XMPP}
{IRC}
{IP Messanger}
{A}
{XMPP$B$O%*!<%W%s$J%$%s%9%?%s%H%a%C%;%s%8%c!<MQ%W%m%H%3%k$N#1$D!#$J$*!"(BDebian$B4X78<T$NMxMQ$K$"$?$C$F>\$7$/$O!"(Bhttp://rtc.debian.org$B$r;2>H!#(B}

\santaku
{2015/10/30$B$K$F!"(BDebian$B$K$F!"(B2$B2sL\$N8xJg$,9T$o$l$?Lr$^$o$j$O<!$N$I$l!)(B}
{2016 DPN}
{technical commitee}
{2016 Debian JP$B2qD9(B}
{B}
{2$BL>$[$I!"(BDebian Developer$B$NJ}$G(Btechnical comittee$B$G3hLv$G$-$kJ}Jg=8$H$N$3$H$G$9!#8=?&(Btechnical comittee$B$N?M$i$,!"(B2015/12/31$B$GG$4|$,@Z$l$F$7$^$&$H$$$&$3$H$G!"$=$l$^$G$K8uJd<T$r5s$2$kI,MW$,$"$k$H$N$3$H$G$9!#(B}



\end{multicols}

\dancersection{最近のDebian関連のミーティング報告}{野島 貴英}

\subsection{第131回東京エリアDebian勉強会}

\begin{itemize}
\item 場所はdotsさんをお借りしての開催でした。
\item 参加者は7名でした。
\item セミナ内容は野島さんによる「DebConf15ビデオ紹介」、Rogerさんによる「毎日使えるIPv6ネットワークの構築」でした。
\item 残りの時間でhack timeを行い、成果発表をしました。
\end{itemize} 

　野島さんによる「DebConf15ビデオ紹介」は、DebConf15の公開ビデオから、
\begin{itemize}
  \item ``Stretching out for trustworthy reproducible builds''
  \item ``Thanks for maintaining a desktop environment. But is it accessible?''
  \item ``Citizenfour Screening''
\end{itemize}
の３つのセッションを選び、サマリとハイライトについて語りました。３つのセッションとも、日本では語られる事が無い珍しい話題となります。是非東京エリアDebian勉強会資料とともに、ビデオを視聴いただくことをおすすめします。

　Rogerさんによる「毎日使えるIPv6ネットワークの構築」は、実際に自宅にDebian機をIPv6環境でつなぐことに関する苦労話や、設定方法についての発表でした。昨今、日本の総務省からも、各携帯3社へ、IPv6接続の関する圧力がかかる状況\footnote{\url{http://www.itmedia.co.jp/news/articles/1511/11/news062.html}}ですので、Debianの熱烈ユーザの皆様におきましては、是非こちらの発表をきっかけにIPv6接続をDebianでお試しください。

\subsection{第132回東京エリアDebian勉強会(OSC 2015 Tokyo/Fall)}

\begin{itemize}
\item 場所は明星大学さんでOSC 2015 Tokyo/Fallのイベントが行われました。
\item 展示用のブースを出し、Debianをこれから初める方、Debianを使って困っている方などいろいろな方から意見、議論をして交流を行いました。
\item セミナの内容は岩松さんによる「Debian と systemd について」でした。
\item ブースには65名、セミナは部屋がほぼ満員でした。
\end{itemize} 


% % (query-replace-regexp "<.*?>" "")
% % (query-replace-regexp "^[	 ]\+" "")

%-------------------------------------------------------------------------------
\dancersection{Debian GNU/kFreeBSD セットアップガイド 2015年版}{杉本 典充}
%-------------------------------------------------------------------------------

\subsection{はじめに}

Debian GNU/kFreeBSDは、FreeBSDカーネルで動作するDebianです。Debianは単なるLinuxディストリビューションではなく、ユニバーサルオペレーティングシステムを目指しており、その例としてDebian GNU/kFreeBSDがあります。\footnote{私が知っている限り「Universal Operating System」という記述は本家webサイトトップページのtitleタグでしか見たことないです。}

ただDebian GNU/kFreeBSDはその特異さゆえにDebian GNU/Linuxと異なる点があります。今回は、Debian GNU/kFreeBSDに触れるにあたり、どのようにセットアップを行うとよいか説明します。

\subsection{Debian PortsとDebian GNU/kFreeBSD}

Debian Ports\footnote{\url{https://www.debian.org/ports/}}とは、さまざまなCPUやカーネルで動作するように移植を行うプロジェクトのことです。

FreeBSDカーネルで起動するDebianを作るプロジェクトがあり、そのDebianのことを「Debian GNU/kFreeBSD」と呼んでいます(kはkernelのこと)。現在ではIntel CPUのアーキテクチャのみあります(kfreebsd-amd64、kfreebsd-i386)。Debian 6 (Squeeze)ではテクノロジープレビューとして初めてstableリリースされました。Debian 7 (Wheezy)でも継続してstableリリースされたのですが、Debian 8 (Jessie)ではDropとなったためリリースされていません。\footnote{\url{https://lists.debian.org/debian-devel-announce/2014/09/msg00002.html} において、stableを維持しつつsidの開発を進めるにはそれなりに人手が必要であるが、kFreeBSDを作業する人手は不足していることが指摘されています。}


\subsection{Debian GNU/kFreeBSDのインストール}
\subsubsection{インストールイメージの入手}

\url{https://www.debian.org/devel/debian-installer/}にあるdaily-imagesを使ってインストールします。
kfreebsd-amd64版のmini.isoをダウンロードし、USBメモリにddしてインストールディスクを作成します。kfreebsd-i386版のmini.isoを利用しても構わないのですが、ファイルシステムにZFSを使う場合はメモリ不足になりがちなため注意してください。

mini.isoをddしたUSBメモリを差してPCを起動するとDebian Installerが起動します。なお、現時点のkfreebsd版Debian Installerは以下の制約があります。\footnote{本家FreeBSD-10.1のインストーラはUEFIブートに対応し、ディスク形式をGPTとしてインストールすることが可能です。}
\begin{itemize}
  \item UEFIブートに対応していない
  \item ディスク形式はMBRのみに対応している(GPTは非対応)
\end{itemize}


\subsubsection{インストーラの表示言語}

kfreebsd版Debian Installerは、日本語の表示ができません(インストーラでフレームバッファが有効になっていないと思われる)。そのため、LANG=Cでインストールを進めます。

\subsubsection{パーティション構成とファイルシステム}

Debian GNU/kFreeBSDをインストールするときはrootパーティションをMBRの基本パーティションにする必要があります（拡張パーティションにインストールするとgrubのインストールに失敗します）。swapパーティションは拡張パーティションに作成しても問題ありません。

この前提があるため、プリインストールのWindowsとデュアルブートしたい場合は、以下のパーティション構成でほぼ決まりになります。\footnote{ファイルシステムをUFSにする場合は基本パーティションの上限(４つ)を超えることがあるため、ディスク内のWindows用リカバリ領域を削除する、Dドライブを削除するなどの事前準備が必要になります。}

\begin{commandline}
# fdisk -l /dev/ada0
Disk /dev/ada0: 298.1 GiB, 320072933376 bytes, 625142448 sectors
Units: sectors of 1 * 512 = 512 bytes
Sector size (logical/physical): 512 bytes / 512 bytes
I/O size (minimum/optimal): 512 bytes / 512 bytes
Disklabel type: dos
Disk identifier: 0x33d61950

Device      Boot     Start       End   Sectors   Size Id Type
/dev/ada0p1 *         2048   2459647   2457600   1.2G  7 HPFS/NTFS/exFAT
  -> Windows 7のインストーラが自動で確保する領域
/dev/ada0p2        2459648 375142399 372682752 177.7G  7 HPFS/NTFS/exFA
  -> Windows 7のOSをインストールしたNTFSパーティション
/dev/ada0p3      580083712 625140399  45056688  21.5G  7 HPFS/NTFS/exFAT
  -> ノートPCメーカーのリカバリー領域
/dev/ada0p4      375142400 580083711 204941312  97.7G a5 FreeBSD
  -> Debian GNU/kFreeBSDをインストールするZFSストレージプール領域
  -> ZFSパーティションの中にrootボリュームとswapボリュームを作成します
\end{commandline}

4つの基本パーティションにうち、1つをZFSストレージプール(LVMの物理ボリュームに似ている概念)に割り当てます。そしてZFSストレージプールからrootボリュームとswapボリュームを作成します(別途bootパーティションを単独で用意しなくても問題ありません)。

ディスクすべてをkfreebsdに割り当てることができ、かつメモリ搭載量が少ない環境(古いPCや仮想マシン)へインストールする場合はUFSを選択したほうがよいと思います。\footnote{インストール時に作成するUFSパーティションはsoft updateが無効になっています(つまり同期書き込みする設定になっている)。rescue modeから"tunefs.ufs -n enable "partition-path""を実行するとsoft updateを有効にして非同期書き込みに変更できますが、耐障害性が低下しますのでご注意ください。}

\subsection{Debian GNU/kFreeBSD固有のDebianパッケージ}

Debian GNU/kFreeBSD固有のパッケージの例を紹介します。

\subsubsection{kfreebsd-imageパッケージ}

Debian GNU/kFreeBSDのkernelイメージを収録したパッケージです。unstableにはkfreebsd-image-10.1があり現状のデフォルトkernelになっています。experimentalにはkfreebsd-image-10.2\footnote{FreeBSD-10.2は2015年8月14日 UTCにリリースされていますので次第にこちらがデフォルトになるでしょう。}、kfreebsd-image-11\footnote{現在のFreeBSD-CURRENTです。FreeBSD-CURRENTは、debianでいうところのunstableと似た位置付けであり、最新の開発版です。}もあります。

\subsubsection{zfsutilsパッケージ}

zfsutilsパッケージはZFSを操作するコマンドを含んだパッケージです。インストール時のファイルシステムにZFSを選択した場合はデフォルトでインストールされます。

kfreebsd-image-10.1で利用できるZFSのバージョンはver 28となっています。

\begin{commandline}
  $ zpool upgrade -v
  (snip)
  28  Multiple vdev replacements
\end{commandline}

\subsubsection{freebsd-utilsパッケージ}

FreeBSD固有のコマンドを含んだパッケージです。/sbin/mount\_*、/usr/sbin/jailなどが入っています。

\subsubsection{freebsd-smbfsパッケージ}

freebsd-smbfsパッケージは、Windowsファイル共有(SMB共有)へアクセスするためのパッケージです。
インストールすると、/usr/sbin/mount\_smbfsコマンドが使えるようになります。

Windowsファイル共有先をmountするには以下のコマンドを実行します。

\begin{commandline}
# mount_smbfs -E UTF-8:CP932 -I {ファイルサーバのIPアドレス} -U {smbユーザ名} //{ファイルサーバのIPアドレス}/{dir} {mount先dir}
\end{commandline}

\subsubsection{freebsd-pppパッケージ}

freebsd-pppパッケージは、ダイアルアップする/usr/sbin/pppコマンドを含んでいます。3GやLTEに対応したUSBモデムを使う場合に必要となります。

Debian GNU/Linuxではpppパッケージやwvdialパッケージでダイアルアップしますが、kfreebsdでは使えないため注意が必要です。

ppp接続の例については後述します。

\subsubsection{pfパッケージ}

FreeBSD kernelがもつPacket Filterと呼ばれるいわゆるファイアウォール機能を制御するコマンド/sbin/pfctlを含んだパッケージです。\footnote{Linux kernelのnetfilter機能と制御コマンドiptablesに相当します。}

/sbin/pfctlの設定ファイルは/etc/pf.confであり、Linuxのiptables用設定ファイルと中身が全く異なります。

\subsection{WindowsとDebian GNU/kFreeBSDのデュアルブート設定}

Debian GNU/kFreeBSDのboot loaderはgrub2を利用しています。
grub2でDebian GNU/kFreeBSDとWindowsをデュアルブートしたい場合は以下の操作を行い、grubにメニューを追加します。((hd0,2)の部分はインストール環境に合わせて変更してください))

\begin{commandline}
  # cd /etc/grub
  # vi 40_custom.kfreebsd
  
  #!/bin/sh
  exec tail -n +3 $0
  # This file provides an easy way to add custom menu entries.  Simply type the
  # menu entries you want to add after this comment.  Be careful not to change
  # the 'exec tail' line above.
  menuentry ``Windows (loader)'' {
    insmod part_msdos
    insmod ntfs
    set root=(hd0,2)
    chainloader +1
  }

  # update-grub
\end{commandline}

\subsection{ハードウェアとソフトウェアのセットアップ}
\subsubsection{有線LAN}

有線LANは利用するドライバによってデバイス名が変化します(intelのPC向けならem0、realtekならre0)。
設定ファイルはDebian GNU/Linuxと同じ/etc/network/interfacesですが、allow-hotplug句はlinuxで使われるudevが提供している機能であることからkfreebsdでは利用できないため注意が必要です。

そのため、有線LAN接続環境がない状況でOSを起動すると有線LANによるDHCPのIPアドレス取得がタイムアウトするまでloginプロンプトが出てこなくなります(起動に時間がかかる)。私はノートPCにDebian GNU/kFreeBSDをインストールする場合は以下コマンドを手動で実行してネットワークへ接続するようにしています。

\begin{commandline}
  # vi /etc/network/interfaces

  #auto em0  <-コメントにします
  iface lan_home inet dhcp

  
  # ifup em0=lan_home
\end{commandline}

\subsubsection{無線LAN}

無線LANはThinkpad X220に搭載している「Intel Centrino advanced-N 6205」で動作することを確認しています。そのため、同じiwnデバイスと認識される「Intel Wireless WiFi Link 4965」以降のIntel製無線LANカードであれば動作すると思います。

無線LANを利用するためにkfreebsd-image-10系の最新版kernel、firmware、無線LANデーモン"wpasupplicant"をインストールします。\footnote{kfreebsdで無線LANが使えるようになったのは2014年8月頃と思われます。 \url{https://bugs.debian.org/cgi-bin/bugreport.cgi?bug=642468}}

\begin{commandline}
  # vi /etc/apt/sources.list

  deb http://ftp.jp.debian.org/debian/ unstable main contrib non-free

  # apt-get install kfreebsd-image-10-amd64
  # apt-get install firmware-iwlwifi wpasupplicant
  # reboot
\end{commandline}

FreeBSDは物理無線LANインタフェースと論理無線LANインタフェースのように分かれています。そのため、論理無線LANインタフェースを生成します。

\begin{commandline}
  # ifconfig iwn0  (これが物理インタフェース名)
  iwn0: flags=8802 metric 0 mtu 2290
  ether xx:xx:xx:xx:xx:xx
  media: IEEE 802.11 Wireless Ethernet autoselect (autoselect)
  status: no carrier
  nd6 options=23

  # ifconfig wlan create wlandev iwn0
  wlan0

  # ifconfig wlan0 (これが論理インタフェース名)
  wlan0: flags=8802 metric 0 mtu 1500
  ether xx:xx:xx:xx:xx:xx
  inet6 fe80::xxxx:xxxx:xxxx:xxxx%wlan0 prefixlen 64 scopeid 0x6
  ssid " channel 5 (2432 MHz 11g)
  country US authmode WPA2/802.11i privacy OFF txpower 15 bmiss 10
  scanvalid 60 bgscan bgscanintvl 300 bgscanidle 250 roam:rssi 7
  roam:rate 5 protmode CTS wme
  media: IEEE 802.11 Wireless Ethernet autoselect (autoselect)
  status: no carrier
      nd6 options=23
\end{commandline}

接続する無線LANアクセスポイントの認証情報設定ファイルを作成します。

\begin{commandline}
  $ wpa_passphrase apname1 appassword > wpa_apname1.conf
  $ cat wpa_apname1.conf
  network={
    ssid=''apname1''
    #psk=''appassword''
    psk=e9fdcb43eba09b6342df30f14275625c8494e534799a82d6639b6124434ea627
  }  
\end{commandline}

無線LANアクセスポイントへ接続し、DHCPでIPアドレスを取得します。IPアドレスは論理インタフェースに付与されます。

\begin{commandline}
  # wpa_supplicant -i wlan0 -c ./wpa_apname1.conf
  Successfully initialized wpa_supplicant
  ioctl[SIOCS80211, op=20, val=0, arg_len=7]: Invalid argument
  ioctl[SIOCS80211, op=20, val=0, arg_len=7]: Invalid argument
  wlan0: Trying to associate with yy:yy:yy:yy:tt:tt (SSID='apname1' freq=2432 MHz)
  wlan0: Associated with yy:yy:yy:yy:yy:yy
  wlan0: WPA: Key negotiation completed with yy:yy:yy:yy:yy:yy [PTK=CCMP GTK=CCMP]
  wlan0: CTRL-EVENT-CONNECTED - Connection to yy:yy:yy:yy:yy:yy completed [id=0 id_str=]
  wlan0: WPA: Group rekeying completed with yy:yy:yy:yy:yy:yy [GTK=CCMP]

  # dhclient wlan0

  # /sbin/ifconfig wlan0
  wlan0: flags=8843 metric 0 mtu 1500
  ether xx:xx:xx:xx:xx:xx
  inet6 fe80::xxxx:xxxx:xxxx:xxxx%wlan0 prefixlen 64 scopeid 0x6
  inet 192.168.1.100 netmask 0xffffff00 broadcast 192.168.1.255
  ssid apname1 channel 5 (2432 MHz 11g) bssid yy:yy:yy:yy:yy:yy
  country US authmode WPA2/802.11i privacy ON deftxkey UNDEF
  AES-CCM 2:128-bit txpower 15 bmiss 10 scanvalid 60 bgscan
  bgscanintvl 300 bgscanidle 250 roam:rssi 7 roam:rate 5 protmode CTS
  wme roaming MANUAL
  media: IEEE 802.11 Wireless Ethernet OFDM/48Mbps mode 11g
  status: associated
          nd6 options=23
\end{commandline}

wpasupplicantコマンドで無線LANアクセスポイントへ接続を試みたがエラーが発生し接続できない場合があります。その場合は以下を試すと接続できる場合があります。

\begin{itemize}
  \item 接続先SSIDを2.4GHz帯のものに変更する。
  \item "\# ifconfig wlan0 -ht40"を実行する。\footnote{デュアルチャネル接続を無効にして、20MHz幅の電波で通信するように指示するコマンドです。}
\end{itemize}


\subsubsection{USBモデムによるダイアルアップ}

3GやLTEのUSBモデムを使ってダイアルアップ接続することができます。NTTドコモから出ている「L-02C」というLTEに対応したUSBモデムを例に接続する手順を説明します。

まずはpppコマンドとUSBモデム処理に使うコマンドをインストールします。

\begin{commandline}
  # apt-get install freebsd-ppp usb-modeswitch 
\end{commandline}

L-02CをPCに差すとCD-ROMデバイスとして認識します。そのため、L-02Cをモデムモードに変更するコマンドを実行します。

\begin{commandline}
  # vi /etc/usb_modeswitch.d/L02C.conf
  ########################################################
  # LG L-02C LTE modem

  DefaultVendor= 0x1004
  DefaultProduct=0x61dd

  TargetVendor= 0x1004
  TargetProduct=0x618f

  MessageContent=5553424312345678000000000000061b000000020000000000000000000000
  NeedResponse=1
  CheckSuccess=20

  # usb_modeswitch -c /etc/usb_modeswitch.d/L02C.conf
\end{commandline}

\begin{commandline}
  # ls /dev/cua*
  /dev/cuaU0.0       /dev/cuaU0.1       /dev/cuaU0.2       /dev/cuaU0.3
  /dev/cuaU0.0.init  /dev/cuaU0.1.init  /dev/cuaU0.2.init  /dev/cuaU0.3.init
  /dev/cuaU0.0.lock  /dev/cuaU0.1.lock  /dev/cuaU0.2.lock  /dev/cuaU0.3.lock
\end{commandline}

次にpppコマンドの設定ファイルを作成し、ppp接続します。利用する回線によって適時APN、ユーザ名、パスワードは変更してください。PPP接続に成功するとtunインタフェースが生成され、IPアドレスが付与されます。

\begin{commandline}
  # vi /etc/ppp/ppp.conf

  default:
    set log Phase Chat LCP IPCP CCP tun command
    ident user-ppp VERSION
    set device /dev/cuaU0.2
    set speed 38400
    set dial ``ABORT BUSY TIMEOUT 2 \
    \"\" \
    AT OK-AT-OK \
    AT+CFUN=1 OK-AT-OK \
    AT+CMEE=2 OK-AT-OK \
    AT+CSQ OK \
    AT+CGDCONT=1,\\\"IP\\\",\\\"apn.ne.jp\\\" OK \
    AT+CGACT? OK-AT-OK \
    AT+CGATT? OK \
    AT+CGCLASS? OK \
    AT+COPS? OK \
    ATD*99***1# CONNECT''
    set timeout 180
    enable dns

  myppp:
    set phone ``*99***1#''
    set authname ``apnuser''
    set authkey ``apnpass''
    set timeout 300
    disable ipv6
    set ifaddr 10.1.0.1/0 10.1.0.2/0 255.255.255.255 0.0.0.0
    add default HISADDR

  # ppp -foreground myppp
\end{commandline}

\subsubsection{ビデオドライバ}

現在、主流のビデオカードはIntel、AMD社のRadeonシリーズ、NVIDIA社のGeForceシリーズがあります。
Debian GNU/kFreeBSDはFreeBSD向けに提供されるプロプラエタリドライバのビルド環境がないため、オープンソース版ドライバを利用する必要があります。そのため、IntelまたはRadeonシリーズのビデオカードを利用することをおすすめします。\footnote{NVIDIA社ビデオカードを利用する場合、xserver-xorg-video-nouveauパッケージはLinux版のみの提供であるため使えず、オープンソース版ドライバxserver-xorg-video-nvは開発がすでに止まっておりかなり古いビデオカードしか対応してません。}

Intel系ビデオカードを利用する場合は以下のパッケージをインストールします。

\begin{commandline}
# apt-get install xserver-xorg-video-intel
\end{commandline}

Radeon系ビデオカードを使用する場合は以下のパッケージをインストールします。

\begin{commandline}
# apt-get install xserver-xorg-video-ati
\end{commandline}

次はKMS(kernel mode settings)を有効にします。以下コマンドを実行するとコンソール画面の解像度が上がります\footnote{昔のkfreebsdでは、i915.koをloadしたままi915kms.koをloadするとrebootしてしまう現象が起こっていたため念のためunloadしています。現在のkfreebsd-image-10.1のkernelで試したところ、loadしたままでも大丈夫ではあるようです。}。KMSを有効にした状態にすると、X Window Systemで液晶モニタの解像度を最大限に利用でき、xrandrコマンドで外部モニタ出力もできるようになります。

\begin{commandline}
  # kldunload i915
  # kldload i915kms
\end{commandline}

再起動後も自動でkernel moduleをロードするように設定します。

\begin{commandline}
  # vi /etc/modules
  i915kms
\end{commandline}

\subsubsection{localeの再設定}

Debian InstallerではLANG=Cを選択してインストールしているため、出力メッセージが英語になっています。そのためlocaleを日本語に変更します。(ただし、コンソール環境では日本語メッセージが化けるので注意)

\begin{commandline}
  # dpkg-reconfigure locales

  -> ja_JP.UTF-8を選択する
\end{commandline}

\subsubsection{X Window Systemのキーボード設定}

FreeBSDのxorgではhalを使ってキーボードのレイアウトを自動判定しています。しかし、halはupstreamによるメンテナンスをすでに終了しており、kfreebsdでもdebianパッケージの提供は終了しています。そのためX Window System起動時のキーボードレイアウトはデフォルトの英語キーボードと判定されます。

キーボードレイアウトの変更はxorg.confで直接指定する、ウィンドウマネージャのキーボード設定を利用するなどして対応する必要があります。

\subsubsection{webブラウザ}

webブラウザはDebian GNU/Linux同様にiceweaselパッケージが提供されています。しかしchromiumパッケージはkfreebsdに存在しません。

またAdobe Flash PlayerはLinux用のバイナリとして提供されるため、Flashを見たい場合はgnashなどをインストールする必要があります(ただし動作の安定度は未知数)。

\subsubsection{USB 3.0}

kfreebsd-image-10-amd64パッケージのkernelにおいて、xhci.koがstatic linkされていることを確認しています。ただし、USB 3.0をもつPCにDebian GNU/kFreeBSDをインストールしたことがないため動作は未確認です。

\subsubsection{サウンドドライバ}

FreeBSDはOSSという仕組みでサウンドを出力します(ALSAはLinux専用のサウンドシステムです)。最近のPCにはHigh Definition Audio規格のチップが搭載されることが多いため、snd\_hda.koドライバでサウンドを出力することができます(snd\_hda.koはkernelにstatic linkされています)。

また、pulseaudioパッケージをインストールし、audaciousを使ってpulseaudio経由で音楽を再生できることは確認しています。

\subsubsection{電源関係}

CPUクロックの制御はpowerdパッケージのpowerdが行っています。現在動作中のCPUクロック数はsysctlコマンドで取得できます。

\begin{commandline}
  $ sysctl dev.cpu.0.freq
  dev.cpu.0.freq: 800
\end{commandline}

バッテリー残量を取得するにはacpiconfコマンドを実行します。

\begin{commandline}
  $ /usr/sbin/acpiconf -i 0
  Design capacity:62160 mWh
  Last full capacity:26300 mWh
  Technology:secondary (rechargeable)
  Design voltage:11100 mV
  Capacity (warn):1315 mWh
  Capacity (low):200 mWh
  (snip)
  State:discharging
  Remaining capacity:95%
  Remaining time:1:25
  Present rate:17681 mW
  Present voltage:12186 mV
\end{commandline}

サスペンドとハイバーネートについては未確認です。

\subsubsection{USBメモリのmount}

FreeBSDではUSBメモリをPCに差すと/dev/da0s1のように認識します。FAT16またはFAT32の領域をもつUSBメモリを/mnt/usbへmountするには以下のコマンドを実行します。

\begin{commandline}
# mount_msdosfs -L ja_JP.UTF-8 -D CP932 /dev/da0s1 /mnt/usb
\end{commandline}

exFATをmountするのはexfat-fuseパッケージを、NTFSをmountするにはntfs-3gパッケージを使うと思われますが、試したところどちらもエラーが出てmountできませんでした。

\subsubsection{Linuxエミュレーション}

FreeBSD kernelにはlinuxバイナリ互換機能があり、Linuxのi386バイナリを実行することができます。\footnote{FreeBSD kernelのkernel moduleであるlinux.koが処理しています}

debootstrapでlinux-i386環境を用意しchrootすることによってlinux-i386形式のバイナリを実行することができます。この機能を使うことによりlinuxバイナリのみ提供されるソフトウェアをkfreebsd上で利用することができます。\footnote{本家Java、Adobe Flash Player、Adobe Readerなどが該当します。}。

\subsubsection{仮想化}

FreeBSDにはFreeBSD Jailというコンテナ型仮想化環境を実行する機能があります。freebsd-utilsパッケージをインストールすることで利用できますが移植が完全ではなく一部機能がない状態です。\footnote{jlsコマンド、jexecコマンドがまだありません。}

Debian GNU/kFreeBSDにqemuパッケージはありますので、他のOSを使う必要がある場合は利用してもよいでしょう。(ただしあまり使ったことがないため、動作の安定度は未知数)

その他にFreeBSDにはOSを完全仮想化して動作させるvirtualbox、bhyveがありますが、Debian GNU/kFreeBSDにはまだ移植されていません。

\subsection{おわりに}

Debian GNU/kFreeBSDのインストール方法とセットアップ方法について説明しました。動作確認ができていない機能や移植されていない機能もまだ多くありますが、OSを開発したい方にはよい腕試しの場になると思います。興味のある方はまず使ってみるところから始めてみるのはいかがでしょうか。

\subsection{参考文献}

\begin{itemize}
  \item 「Debian\_GNU/kFreeBSD - Debian Wiki」 \url{https://wiki.debian.org/Debian\_GNU/kFreeBSD}
  \item 杉本典充 (2015). 「Debian GNU/kFreeBSDにおけるJail構築を試してみた」 \url{http://tokyodebian.alioth.debian.org/pdf/debianmeetingresume201502-presentation-sugimoto.pdf}
\end{itemize}

%-------------------------------------------------------------------------------
%-------------------------------------------------------------------------------
\dancersection{会場での無線LANのつなぎ方}{野島 貴英,Roger}
%-------------------------------------------------------------------------------
 \subsection{はじめに}

　今回会場側には無線LAN経由のグローバル回線が用意されています。

　以下にDebianマシンでの接続方法を記載します。

 また、自分の環境では違うやり方でつながったという方は、野島まで
教えて下さい。こちらでもノウハウとして溜めていく予定です。

 \subsection{wpasupplicant及び/etc/network/interfacesを利用の場合}

 もっとも良いマニュアルは、/usr/share/doc/wpasupplicant/README.Debian.gz
となります。困った場合はこちらも合わせてご参照下さい。

　以下に/etc/network/interfacesの定義について会場の例を記載します。

\begin{commandline}
$ sudo vi /etc/network/interfaces
-----以下のエントリがなければ追記ここから----------
iface wlan0_debian inet dhcp
     wpa-conf /etc/wpa_supplicant/wpa_supplicant_debian.conf
-----以下のエントリがなければ追記ここまで----------
$ sudo vi /etc/wpa_supplicant/wpa_supplicant_debian.conf
-----以下のエントリを追記ここから----------
network={
     ssid=<<会場のSSID>>
     psk=<<会場のパスワード>>
     scan_ssid=1
}
-----以下のエントリを追記ここまで----------
$ sudo chmod 600 /etc/wpa_supplicant/wpa_supplicant_debian.conf
$ sudo ifup wlan0=wlan0_debian
\end{commandline}
%$

 また、ハマってしまった時のデバッグ方法は、
/usr/share/doc/wpasupplicant/README.Debian.gz中の''4. Trubleshooting''の章が便利です。

 \subsection{その他の無線LAN用パッケージを利用の場合}

　すみません、自分が情報を持たないため、現場で教えて下さい。

\cleartooddpage

\vspace*{15cm}
\hrule
\vspace{2mm}
\includegraphics[width=2cm]{image200502/openlogo-nd.eps}
\noindent \Large \bf Debian 勉強会資料\\
\noindent \normalfont \debmtgyear{}年\debmtgmonth{}月\debmtgdate{}日 \hspace{5mm}  初版第1刷発行\\
\noindent \normalfont 東京エリア Debian 勉強会 （編集・印刷・発行）\\
\hrule

\end{document}
